\documentclass{article}
\usepackage{enumitem}
\usepackage{amsmath}
\begin{document}

\title{Chapter 34: Biotechnology and its Applications}
\date{}
\maketitle

\begin{enumerate}[label=Q\arabic*.]

\item Transgenic \textit{Bt} crops produce insecticidal Cry proteins. These proteins:
\\(A) Repel insects by smell
\\(B) Are toxic only to specific insects (e.g., lepidopteran larvae) by forming pores in the midgut at alkaline pH; harmless to mammals
\\(C) Kill all insects including beneficial ones like bees equally
\\(D) Inhibit photosynthesis in insects

\item RNA interference (RNAi) was used to create a nematode-resistant transgenic plant by:
\\(A) Expressing nematode Cry protein in the plant
\\(B) Using dsRNA targeting a specific nematode gene, silencing it when the nematode feeds on the plant
\\(C) Introducing a nematicide gene into the plant
\\(D) Over-expressing anti-fungal proteins

\item Gene therapy involves:
\\(A) Treatment of disease by administering drugs that target genes
\\(B) Introduction of a correct copy of a gene into the cells of a patient with a genetic disorder
\\(C) Editing genes in germline cells only
\\(D) Producing therapeutic proteins outside the body

\item ADA (adenosine deaminase) deficiency was the first disease treated by gene therapy. The approach involved:
\\(A) Bone marrow transplant only
\\(B) Introducing functional ADA gene into patient's T lymphocytes ex vivo using a retroviral vector
\\(C) Enzyme replacement therapy without genetic modification
\\(D) CRISPR-Cas9 editing of the patient's genome directly

\item Assertion: Transgenic animals are valuable models for studying human diseases.\\
Reason: Transgenic animals express specific human genes or mutations, mimicking human disease pathology.
\\(A) Both true, Reason is correct explanation
\\(B) Both true, Reason is not correct explanation
\\(C) Assertion true, Reason false
\\(D) Assertion false, Reason true

\item Rosie, the first transgenic cow, was significant because she:
\\(A) Produced human growth hormone in milk
\\(B) Produced human protein-enriched milk containing human alpha-1-antitrypsin
\\(C) Was cloned from a somatic cell
\\(D) Was resistant to bovine tuberculosis

\item Molecular diagnostics using PCR can detect pathogens:
\\(A) Only after culture of the pathogen
\\(B) Even in very small quantities from a patient sample before symptoms appear (early detection)
\\(C) Only from blood samples
\\(D) Only after immune response has developed

\item The herbicide-resistant crop (e.g., Roundup Ready soybean) was created by:
\\(A) Mutating the gene encoding the herbicide target enzyme
\\(B) Introducing a modified EPSPS gene from bacteria that is not inhibited by glyphosate (Roundup)
\\(C) Expressing a herbicide-degrading enzyme
\\(D) Reducing herbicide uptake by modifying roots

\item Bioremediation using genetically engineered microorganisms involves:
\\(A) Using microbes to produce biofuels
\\(B) Using engineered microbes to degrade environmental pollutants (e.g., oil spills, heavy metals)
\\(C) Producing antibiotics from soil microbes
\\(D) Using microbes to fix nitrogen in polluted soils

\item The ``Flavr Savr'' tomato was engineered to:
\\(A) Produce higher lycopene content
\\(B) Delay ripening by antisense suppression of polygalacturonase (PG) gene
\\(C) Resist insect pests
\\(D) Tolerate drought

\item Pharming refers to:
\\(A) Pharmaceutical manufacturing in fermenters
\\(B) Production of pharmaceutical proteins (e.g., antibodies, hormones) in transgenic animals or plants
\\(C) Organic farming for drug plant production
\\(D) Farming using automated systems

\item CRISPR-Cas9 genome editing works by:
\\(A) Randomly mutagenizing the genome
\\(B) Guide RNA directing Cas9 endonuclease to a specific DNA sequence for precise cutting and editing
\\(C) Inserting genes randomly into the genome
\\(D) Silencing genes using antisense RNA only

\item Ethical concerns about GM crops include all EXCEPT:
\\(A) Potential harm to non-target organisms
\\(B) Reduction in crop yield
\\(C) Development of ``superweeds'' through gene flow
\\(D) Biopiracy and corporate control of seed supply

\item Match the following transgenic organisms with their applications:\\
1. \textit{Bt} cotton $\rightarrow$ (a) Pest resistance (bollworm)\\
2. Golden Rice $\rightarrow$ (b) Provitamin A production\\
3. Rosie (cow) $\rightarrow$ (c) Human alpha-1-antitrypsin in milk\\
4. \textit{E. coli} (recombinant) $\rightarrow$ (d) Human insulin production
\\(A) 1-a, 2-b, 3-c, 4-d
\\(B) 1-b, 2-a, 3-d, 4-c
\\(C) 1-c, 2-d, 3-a, 4-b
\\(D) 1-d, 2-c, 3-b, 4-a

\item Biopiracy refers to:
\\(A) Illegal cloning of organisms
\\(B) Unauthorized exploitation of biological resources and traditional knowledge of developing countries by corporations without fair compensation
\\(C) Theft of laboratory equipment
\\(D) Illegal trade in endangered species

\end{enumerate}
\end{document}
